%============================================================
\chapter{Trách Nhiệm Và Danh Dự (Responsibility and Honor)}
\begin{center}
  \textit{\color{truyenblue}The price of greatness is responsibility.}\\
  \textit{(Giá của sự vĩ đại là trách nhiệm.)}\\[4pt]
  \small\color{truyengold}--- Winston Churchill ---
\end{center}
\ngancach

\section{Nguyễn Trãi -- Gánh Nặng Trách Nhiệm Với Dân Tộc}
\begin{truyen}{Nguyễn Trãi -- Mười Năm Nằm Gai Nếm Mật}{Việt Nam, Thế Kỷ XV}
\chuhoa{N}{}N guyễn Trãi chứng kiến cha bị giặc Minh bắt đem về Trung Hoa khi ông còn trẻ.\\
\textit{(Nguyen Trai witnessed his father being taken captive to China by the Ming invaders when he was still young.)}

Cha dặn: ``Con không được theo ta. Con phải ở lại gánh vác trách nhiệm với đất nước.''\\
\textit{(His father instructed: ``Do not follow me. You must stay and shoulder responsibility for our country.'')}

Nguyễn Trãi chôn nỗi đau riêng, tìm đến Lê Lợi dâng ``Bình Ngô Sách'' -- kế sách đánh giặc.\\
\textit{(Nguyen Trai buried his personal grief, sought out Le Loi and presented the ``Binh Ngo Sach'' -- strategy to defeat the invaders.)}

Suốt mười năm kháng chiến, ông là trí não của cuộc khởi nghĩa -- viết hịch, thảo văn thư, lập kế mưu.\\
\textit{(For ten years of resistance, he was the brain of the uprising -- writing proclamations, drafting letters, devising strategies.)}

``Bình Ngô Đại Cáo'' -- bản tuyên ngôn độc lập đầu tiên của Việt Nam -- là tinh hoa của trách nhiệm với dân tộc.\\
\textit{(The ``Binh Ngo Dai Cao'' -- Vietnam's first declaration of independence -- was the finest fruit of his responsibility to his people.)}

\end{truyen}
\begin{baihoc}
  Trách nhiệm với cộng đồng đôi khi đòi hỏi ta gác lại nỗi đau riêng để phục vụ điều lớn lao hơn.\\
  \textit{(Responsibility to community sometimes requires us to set aside personal pain to serve something greater.)}
\end{baihoc}

\section{Harry Truman Và Tấm Biển ``The Buck Stops Here''}
\begin{truyen}{Harry Truman -- Trách Nhiệm Dừng Lại Ở Đây}{Nước Mỹ, Thế Kỷ XX}
\chuhoa{T}{}T ổng thống Harry Truman có một tấm biển gỗ nhỏ để trên bàn làm việc: ``The Buck Stops Here.''\\
\textit{(President Harry Truman kept a small wooden sign on his desk: ``The Buck Stops Here.'')}

Câu đó có nghĩa: mọi trách nhiệm cuối cùng đều dừng tại bàn này -- không đùn đẩy lên trên.\\
\textit{(The phrase meant: all final responsibility stops at this desk -- no passing the buck upward.)}

Khi đưa ra quyết định khó nhất lịch sử -- thả bom nguyên tử -- ông không đổ lỗi cho ai.\\
\textit{(When making the most difficult decision in history -- dropping the atomic bomb -- he blamed no one.)}

Ông nói: ``Tôi là người ra quyết định. Tôi là người chịu trách nhiệm. Lịch sử sẽ phán xét tôi.''\\
\textit{(He said: ``I am the decision-maker. I bear the responsibility. History will judge me.'')}

Trước khi rời Nhà Trắng, ông viết: ``Nếu bạn không chịu được áp lực thì đừng vào bếp nấu.''\\
\textit{(Before leaving the White House, he wrote: ``If you can't stand the heat, stay out of the kitchen.'')}

\end{truyen}
\begin{baihoc}
  Lãnh đạo thực sự là nhận trách nhiệm, không tìm cách đổ lỗi. Trách nhiệm là giá của quyền lực và danh dự.\\
  \textit{(True leadership is accepting responsibility, not finding someone to blame. Responsibility is the price of power and honor.)}
\end{baihoc}

\section{Đặng Lê Nguyên Vũ Và Trách Nhiệm Với Cà Phê Việt}
\begin{truyen}{Đặng Lê Nguyên Vũ -- Giấc Mơ Cà Phê Việt Ra Thế Giới}{Việt Nam Hiện Đại}
\chuhoa{Đ}{}Đ ặng Lê Nguyên Vũ xuất thân nghèo, lớn lên ở vùng quê cà phê Tây Nguyên.\\
\textit{(Dang Le Nguyen Vu came from poverty, growing up in the coffee-growing Central Highlands.)}

Ông chứng kiến nông dân trồng cà phê ngon nhất thế giới nhưng chỉ bán thô với giá rẻ mạt.\\
\textit{(He witnessed farmers growing some of the world's finest coffee but selling it raw at rock-bottom prices.)}

Ông tự nhận trách nhiệm: ``Ai đó phải xây dựng thương hiệu cà phê Việt tầm quốc tế -- tại sao không phải là tôi?''\\
\textit{(He took on the responsibility himself: ``Someone must build an internationally recognized Vietnamese coffee brand -- why not me?'')}

Từ một chiếc máy xay tay và vài ký cà phê, Trung Nguyên lớn lên thành thương hiệu toàn cầu.\\
\textit{(From a hand grinder and a few kilos of coffee, Trung Nguyen grew into a global brand.)}

Ông nói: ``Trách nhiệm không ai giao cho tôi -- tôi tự nhận lấy vì tôi thấy điều đó cần phải làm.''\\
\textit{(He said: ``No one assigned me this responsibility -- I took it myself because I saw it needed to be done.'')}

\end{truyen}
\begin{baihoc}
  Danh dự lớn nhất đến từ việc tự nhận trách nhiệm với cộng đồng mình mà không đợi ai giao phó.\\
  \textit{(The greatest honor comes from voluntarily taking responsibility for your community without waiting for anyone to assign it.)}
\end{baihoc}

\section{Leonidas Và Ba Trăm Người Sparta}
\begin{truyen}{Leonidas -- Danh Dự Của Người Chiến Binh}{Hy Lạp Cổ Đại, 480 TCN}
\chuhoa{V}{}V ua Leonidas dẫn 300 chiến binh Sparta và 700 người đồng minh chặn 300.000 quân Ba Tư tại hẻm núi Thermopylae.\\
\textit{(King Leonidas led 300 Spartan warriors and 700 allies to block 300,000 Persian troops at the pass of Thermopylae.)}

Biết mình sẽ chết, ông vẫn không rút lui vì đã thề bảo vệ Hy Lạp.\\
\textit{(Knowing he would die, he still did not retreat because he had sworn to protect Greece.)}

Khi sứ giả Ba Tư yêu cầu nộp khí giới, Leonidas đáp: ``Molon labe'' -- ``Hãy tự đến mà lấy.''\\
\textit{(When the Persian envoy demanded they surrender their weapons, Leonidas replied: ``Molon labe'' -- ``Come and take them.'')}

Ba ngày cầm chân đại quân Ba Tư đã cho Hy Lạp đủ thời gian chuẩn bị -- thay đổi kết cục lịch sử.\\
\textit{(Three days holding back the Persian army gave Greece enough time to prepare -- changing the course of history.)}

Thermopylae trở thành biểu tượng vĩnh cửu của danh dự: chết đúng vị trí còn hơn sống mà đào ngũ.\\
\textit{(Thermopylae became an eternal symbol of honor: dying at your post is better than living as a deserter.)}

\end{truyen}
\begin{baihoc}
  Danh dự không nằm ở chiến thắng mà nằm ở sự trung thành với cam kết -- dù phải trả giá bằng tất cả.\\
  \textit{(Honor does not lie in victory but in faithfulness to one's commitment -- even at the cost of everything.)}
\end{baihoc}

\section{Cuộc Hành Quân Của Ernie Pyle}
\begin{truyen}{Ernie Pyle -- Nhà Báo Gánh Trách Nhiệm Với Lính}{Chiến Tranh Thế Giới II}
\chuhoa{E}{}E rnie Pyle là phóng viên chiến trường nổi tiếng nhất Thế Chiến II -- nhưng không theo kiểu anh hùng an toàn từ xa.\\
\textit{(Ernie Pyle was the most famous war correspondent of WWII -- but not in the comfortable safe way of heroes from afar.)}

Ông ăn ngủ cùng lính, chịu đựng bùn lầy và đạn bom như mọi người lính thường.\\
\textit{(He ate and slept with the soldiers, enduring mud and shrapnel like every ordinary fighting man.)}

Ông không viết về chiến thắng hay tướng lĩnh -- ông viết về từng người lính với tên đầy đủ và quê hương cụ thể.\\
\textit{(He did not write about victories or generals -- he wrote about each soldier by full name with their specific hometown.)}

``Đây là cách duy nhất để gia đình họ biết con họ đang làm gì. Đó là trách nhiệm của tôi.''\\
\textit{(``This is the only way their families can know what their sons are doing. That is my responsibility.'')}

Ông chết trên chiến trường Okinawa năm 1945 -- đúng vị trí mình đã chọn.\\
\textit{(He died on the battlefield of Okinawa in 1945 -- in the very place he had chosen to be.)}

\end{truyen}
\begin{baihoc}
  Trách nhiệm thực sự đòi hỏi sự hiện diện -- không thể thực hiện từ khoảng cách an toàn. Danh dự là sống đúng với điều mình đã cam kết.\\
  \textit{(True responsibility demands presence -- it cannot be fulfilled from a safe distance. Honor is living true to what you have committed.)}
\end{baihoc}
